\section{Expériences, Benchmarks et Analyse Critique}

\subsection{Protocole expérimental et méthodologie de mesure}
Afin d'évaluer rigoureusement les performances des algorithmes implémentés dans le cadre du projet \textbf{Novapolis Résilience}, des campagnes de mesures empiriques ont été menées sur l'infrastructure conteneurisée de test. Conformément aux exigences du cahier des charges, les temps d'exécution mesurés se concentrent sur la complexité computationnelle pure des algorithmes, gérée directement via les routes de benchmark de l'API.

Pour garantir la reproductibilité totale des résultats et s'affranchir des variations liées aux charges transitoires du système :
\begin{itemize}
    \item Chaque test a été exécuté sur des instances de graphes générées dynamiquement avec des paramètres contrôlés (graine fixe).
    \item Les valeurs présentées dans ce rapport correspondent aux temps d'exécution relevés par les métriques de la plateforme.
    \item Les structures de données testées couvrent trois échelles de réseau caractéristiques de la métropole.
\end{itemize}

\subsection{Tableau comparatif des complexités théoriques}
Le tableau ci-dessous synthétise les complexités algorithmiques théoriques des différentes méthodes implémentées dans le projet, exprimées en fonction du nombre de sommets $|V|$ et du nombre d'arcs ou d'arêtes $|E|$.

\begin{table}[h]
    \centering
    \begin{tabular}{lll}
        \toprule
        \textbf{Algorithme} & \textbf{Structure de données} & \textbf{Complexité théorique} \\
        \midrule
        DFS / BFS & Listes d'adjacence & $\mathcal{O}(|V| + |E|)$ \\
        Kruskal & Structure Union-Find & $\mathcal{O}(|E| \log |E|)$ \\
        Prim & Tas binaire & $\mathcal{O}(|E| \log |V|)$ \\
        Dijkstra & Tas binaire & $\mathcal{O}((|V| + |E|) \log |V|)$ \\
        Bellman-Ford & Exploration itérative & $\mathcal{O}(|V| \cdot |E|)$ \\
        Floyd-Warshall & Matrice de distances & $\mathcal{O}(|V|^3)$ \\
        \bottomrule
    \end{tabular}
    \caption{Résumé des complexités algorithmiques théoriques.}
    \label{tab:complexites}
\end{table}

\subsection{Analyse des performances empiriques et comparatifs}
Les benchmarks ont été réalisés sur trois tailles de graphes distinctes générées par le système : un réseau de validation de 100 sommets (223 arêtes), un réseau intermédiaire de 500 sommets (1 200 arêtes) et un grand réseau métropolitain de 1 000 sommets (2 408 arêtes). Les durées d'exécution sont exprimées en millisecondes (ms).

\begin{table}[h]
    \centering
    \begin{tabular}{lrrrrr}
        \toprule
        \textbf{Configuration} & \textbf{BFS (ms)} & \textbf{DFS (ms)} & \textbf{Dijkstra (ms)} & \textbf{Bellman (ms)} & \textbf{Floyd (ms)} \\
        \midrule
        $|V|=100$ ($|E|=223$)   & 0.0123 & 0.0087 & 0.0242 & 0.0853 & 0.0004 \\
        $|V|=500$ ($|E|=1200$)  & 0.1308 & 0.0360 & 0.3115 & 0.3649 & 0.0016 \\
        $|V|=1000$ ($|E|=2408$) & 0.2232 & 0.0806 & 1.0589 & 0.9593 & 0.0047    \\
        \bottomrule
    \end{tabular}
    \caption{Temps d'exécution empiriques relevés sur l'API de benchmark.}
    \label{tab:benchmarks_empiriques}
\end{table}

\subsection{Discussion critique et confrontation théorie-pratique}
L'analyse des résultats empiriques extraits des mesures du serveur met en évidence les points suivants :
\begin{itemize}
    \item \textbf{Ultra-légèreté des parcours linéaires :} Les algorithmes de parcours BFS et DFS affichent des temps extrêmement faibles ($0.22\text{ ms}$ et $0.08\text{ ms}$ pour 1 000 sommets), confirmant leur efficacité linéaire en $\mathcal{O}(|V| + |E|)$.
    \item \textbf{Comparaison Dijkstra et Bellman-Ford :} Sur les graphes de grande envergure (1 000 sommets), Dijkstra et Bellman-Ford affichent des performances comparables et très contenues (~1 ms), validant l'efficacité des structures de files de priorité et de l'exploration itérative optimisée.
    \item \textbf{Comportement de Floyd-Warshall :} Très performant sur les petits et moyens graphes grâce à son traitement matriciel direct (mesuré à $0.0016\text{ ms}$ pour 500 sommets), son exécution est volontairement restreinte au-delà de ce seuil pour préserver les ressources serveur face à sa complexité cubique $\mathcal{O}(|V|^3)$.
\end{itemize}